\subsubsection*{5.2. 计算的思路.}

从现在起假设 $\check X=\varphi(\mathbb P(N))$ 的维数为 $r-1$, 从而它是 $\check{\mathbb P}^{r}$ 中的超曲面. 它的次数自然可用 Chern 类 [1] 计算.

对每个光滑射影簇 $T$, 以 $\operatorname{CH}(T)$ 表示其按余维分次的 Chow 环. 对 $z\in\operatorname{CH}(T)$, 以
\[
\int_T z
\]
表示其零维分量的次数. 只要不致混淆(例如在积分号下), $T$ 上可逆层 $\mathcal L$ 在 $\operatorname{CH}^1(T)$ 中的类也记作 $\mathcal L$.

设 $L$ 为 $\operatorname{CH}(\check{\mathbb P}^{r})$ 中一条直线的类. 则
\begin{equation}\tag{5.2.1}
\deg(\check X)
=\int_{\check{\mathbb P}^{r}}L\cdot\check X.
\end{equation}
就集合而言, $\check X$ 为 $\mathbb P(N)$ 在 $\varphi$ 下的像(3.1.2). 但从固有推出
\[
\varphi_*:\operatorname{CH}(\mathbb P(N))
\longrightarrow\operatorname{CH}(\check{\mathbb P}^{r}),
\]
的角度看, 有
\[
\varphi_*[\mathbb P(N)]=\deg(\varphi)[\check X].
\]
因此 (5.2.1) 化为
\begin{equation}\tag{5.2.2}
\deg(\check X)
=\frac{1}{\deg(\varphi)}
\int_{\check{\mathbb P}^{r}}L\cdot\varphi_*[\mathbb P(N)].
\end{equation}
此处滥用记号, $\deg(\varphi)$ 中的 $\varphi$ 表示由 (3.1.2) 诱导的态射 $\mathbb P(N)\to\check X$; 其次数为非负整数, 约定它恰在 $\dim\check X<r-1=\dim\mathbb P(N)$ 时为零.
\footnote{印本写为 $\dim X<r-1$. 到对偶簇的诱导映射一般有限的条件为 $\dim\check X=r-1$.}
投影公式将 (5.2.2) 改写为
\begin{equation}\tag{5.2.3}
\deg(\check X)
=\frac{1}{\deg(\varphi)}
\int_{\mathbb P(N)}\varphi^*(L).
\end{equation}

设 $\eta\in\operatorname{CH}(\check{\mathbb P}^{r})$ 为一个超平面的类. 由于 $L=\eta^{r-1}$, 等式 (5.2.3) 化为
\begin{equation}\tag{5.2.4}
\deg(\check X)
=\frac{1}{\deg(\varphi)}
\int_{\mathbb P(N)}\bigl(\varphi^*(\eta)\bigr)^{r-1}.
\end{equation}
现在具体计算 $\varphi^*(\eta)$.

以 $\xi$(分别为 $\tau$)表示 $\mathbb P(N)$ 上的典范可逆层 $\mathcal O_{\mathbb P(N)}(1)$(分别为 $\mathbb P(\mathcal O_X(1)^{r+1})$ 上的 $\mathcal O_{\mathbb P(\mathcal O_X(1)^{r+1})}(1)$). 由函子性及 (5.1.6), 有
\begin{equation}\tag{5.2.5}
\alpha^*(\tau)=\xi.
\end{equation}
在等同
\[
\mathbb P\bigl(\mathcal O_X(1)^{r+1}\bigr)
\simeq X\times\check{\mathbb P}^{r},
\]
下, 记 $H=\mathcal O_X(1)$, 显然有
\begin{equation}\tag{5.2.6}
\tau=\operatorname{pr}_1^*(H)\otimes\operatorname{pr}_2^*(\eta).
\end{equation}
设 $\pi:\mathbb P(N)\to X$ 为结构态射. 则
\[
\pi=\operatorname{pr}_1\circ\alpha,
\qquad
\varphi=\operatorname{pr}_2\circ\alpha.
\]
将 $\alpha^*$ 作用于 (5.2.6), 并使用 (5.2.5), 得到
\begin{equation}\tag{5.2.6$'$}
\xi=\alpha^*(\tau)=\pi^*(H)\otimes\varphi^*(\eta).
\end{equation}
\footnote{印本将此公式与前一个公式都编号为 (5.2.6); 撇号用于区分第二次出现.}
因此在 $\operatorname{CH}(\mathbb P(N))$ 中,
\begin{equation}\tag{5.2.7}
\varphi^*(\eta)=\xi-\pi^*(H),
\end{equation}
而 (5.2.3) 化为
\begin{equation}\tag{5.2.8}
\deg(\check X)
=\frac{1}{\deg(\varphi)}
\int_{\mathbb P(N)}\bigl(\xi-\pi^*(H)\bigr)^{r-1}.
\end{equation}
(5.2.8) 的右边用以下引理计算(参见 5.3.3).